% ============================================================
% Préambule LaTeX — Rapport Ossatures de bâtiments
% CHEBAP/TS – Fatoumata Bineta SALL – 2025-2026
% ============================================================

\usepackage[utf8]{inputenc}
\usepackage[T1]{fontenc}
\usepackage[french]{babel}
\usepackage[top=2.2cm,bottom=2.5cm,left=2.5cm,right=2.5cm]{geometry}

\usepackage{amsmath,amssymb,mathtools}

\usepackage[dvipsnames]{xcolor}
\definecolor{bleu}{RGB}{30,80,125}
\definecolor{gris}{RGB}{110,110,110}
% Couleurs du logo CHEBAP (échantillonnées sur l'énoncé)
\definecolor{chebaprouge}{RGB}{212,42,10}    % CHE en rouge vermillon
\definecolor{chebapgris}{RGB}{126,123,122}   % BAP en gris moyen

\usepackage{tcolorbox}
\tcbuselibrary{skins,breakable}
\usepackage{tikz}
\usepackage{float}
\usepackage{graphicx}

% Note : \graphicspath sera défini dans main.tex pour pointer vers tous les
% sous-dossiers images/ des séances.

% Pour intégrer le scan de l'énoncé en début de rapport
\usepackage{pdfpages}

% Résultat encadré
\newcommand{\encadre}[1]{%
  \begin{tcolorbox}[
    enhanced, colback=white, colframe=black,
    boxrule=0.8pt, arc=0pt, sharp corners,
    halign=center,
    left=6pt, right=6pt, top=4pt, bottom=4pt
  ]
  #1
  \end{tcolorbox}\vspace{2pt}%
}

% Espace réservé pour image — Renommer ton image en "<nom>.png" puis la mettre
% dans le dossier "images/" de la séance correspondante
\newcommand{\figureplaceholder}[2]{%
  \begin{figure}[H]
    \centering
    \fbox{\begin{minipage}[c][5cm][c]{0.7\linewidth}
      \centering
      \color{gris}\itshape\small
      \textbf{[ Image à insérer ]}\\[4pt]
      Renommer votre fichier en :\\
      \texttt{#1.png}\\[4pt]
      Légende : #2
    \end{minipage}}
    \caption{#2}
    \label{fig:#1}
  \end{figure}%
}
% Pour utiliser l'image après l'avoir mise dans le dossier images/, remplacer par :
% \begin{figure}[H]\centering
%   \includegraphics[width=0.7\linewidth]{nom_image.png}
%   \caption{...}\label{fig:nom_image}
% \end{figure}

\usepackage{setspace}
\setstretch{1.5}
\setlength{\parindent}{1.2cm}
\setlength{\parskip}{3pt}

\usepackage{booktabs,array,tabularx}
\usepackage{multirow}
\renewcommand{\arraystretch}{1.25}

\usepackage{enumitem}
\setlist[itemize]{label=---,left=0.6cm,topsep=2pt,itemsep=1pt}

\usepackage{titlesec}
\titleformat{\section}{\large\bfseries\color{bleu}}{}{0em}{}[{\color{bleu}\vspace{-4pt}\rule{\linewidth}{0.5pt}}]
\titleformat{\subsection}{\normalsize\bfseries\color{bleu}}{}{0em}{}
\titleformat{\subsubsection}{\normalsize\bfseries}{}{0em}{}
\titlespacing{\section}{0pt}{18pt}{8pt}
\titlespacing{\subsection}{0pt}{12pt}{4pt}
\titlespacing{\subsubsection}{0pt}{8pt}{3pt}

\setcounter{secnumdepth}{3}
\setcounter{tocdepth}{2}

\usepackage{hyperref}
\hypersetup{colorlinks,linkcolor=bleu,urlcolor=bleu}

% Bandeau de séance : page de garde dédiée + entrée TOC en niveau 1.
%
% Usage :
%   \seance{N}{titre}        -> sous-sections numérotées N.1, N.2, ...
%   \seance[N]{N \& M}{titre} -> idem, avec N en argument optionnel
%
% L'argument optionnel sert pour les séances groupées (« 5 \& 6 »,
% « 11 \& 12 ») où le premier argument n'est pas un nombre exploitable
% par \setcounter. Par défaut, on suppose que #1 est un entier et on
% l'utilise directement.
\newcommand{\seance}[3][]{%
  \clearpage
  \def\seancenum{\ifx\\#1\\#2\else#1\fi}%
  \setcounter{section}{\numexpr\seancenum-1\relax}%
  \refstepcounter{section}%
  \phantomsection
  \addcontentsline{toc}{section}{Séance #2 : #3}%
  \thispagestyle{empty}
  \vspace*{\fill}
  \begin{tcolorbox}[
    enhanced, colback=bleu, colframe=bleu,
    boxrule=0pt, arc=0pt, sharp corners,
    fontupper=\color{white}\Huge\bfseries,
    halign=center,
    left=20pt, right=20pt, top=30pt, bottom=30pt
  ]
  Séance #2\\[12pt]{\LARGE\mdseries #3}
  \end{tcolorbox}
  \vspace*{\fill}
  \clearpage
}

% Variante : bandeau de section spéciale (énoncé, synthèse) sans
% numéro de séance, mais avec le même style graphique. Ces sections
% n'ont pas de \subsection numérotée à l'intérieur, donc on
% n'incrémente pas le compteur de section : la numérotation des
% séances suit ainsi 1, 2, ..., 10 sans décalage.
% Usage : \sectionspeciale{Énoncé du projet}
\newcommand{\sectionspeciale}[1]{%
  \clearpage
  \phantomsection
  \addcontentsline{toc}{section}{#1}%
  \thispagestyle{empty}
  \vspace*{\fill}
  \begin{tcolorbox}[
    enhanced, colback=bleu, colframe=bleu,
    boxrule=0pt, arc=0pt, sharp corners,
    fontupper=\color{white}\Huge\bfseries,
    halign=center,
    left=20pt, right=20pt, top=30pt, bottom=30pt
  ]
  #1
  \end{tcolorbox}
  \vspace*{\fill}
  \clearpage
}
